%! TeX root = master.tex
% !TEX program = lualatex

\lecture{14}{26-05-2026}{Review}{}

\chapter{Review}
This chapter ties together the whole course. The point is not to re-derive
anything, but to see the \textbf{connections}: most of the supervised methods
share a single linear model and differ only by their loss, kernelization is the
same trick applied to many algorithms, and deep learning is the idea of learning
the representation instead of designing it. The same backbone reappears
everywhere.

\section{The machine learning pipeline}
Machine learning = Data \textrightarrow{} Algorithms \textrightarrow{} Insight.
\begin{itemize}
  \item \textbf{Data} can be tabular attributes, text, speech/sound or images.
  In all algorithms seen, a sample is a vector $x_i \in \mathbb{R}^D$.
  \item \textbf{Insight} can be a category, a description, a vector-valued
  quantity or a different (lower-dimensional) representation.
  \item \textbf{Algorithms} split into \textbf{supervised learning} (relies on
  annotated data, the annotations being the desired insight) and
  \textbf{unsupervised learning} (analyzes or transforms unlabeled data).
\end{itemize}

\section{Supervised learning}
\subsection{Setting}
Supervised learning has two stages: \textbf{training}, where data with
ground-truth labels is used to optimize the model parameters by minimizing a
loss, and \textbf{testing}, where the fixed model predicts the output of a new
sample. Two problem types:
\begin{itemize}
  \item \textbf{Classification}: predict one discrete (unordered) label.
  \item \textbf{Regression}: predict continuous (ordered) value(s).
\end{itemize}

\subsection{The linear model}
The output varies linearly with the inputs. Absorbing the bias by appending a
$1$ to $x$ (so $x \in \mathbb{R}^{D+1}$), a single-output model is
\[
  y = w^{(0)} + \sum_{j=1}^{D} w^{(j)} x^{(j)} = w^T x .
\]
For $C$ output dimensions we use a matrix $W \in \mathbb{R}^{(D+1)\times C}$ and
\[
  \hat{y} = W^T x ,
\]
where each column $w_{(k)}$ predicts one output dimension.

\subsection{Linear regression}
``Best fit'' is encoded as minimizing the squared residuals:
\[
  \min_{w} \sum_{i=1}^{N} \left(x_i^T w - y_i\right)^2 .
\]
This is convex and has a \textbf{closed-form} solution (zero the gradient).
Prediction for a new sample is $\hat{y} = (w^*)^T x$.

\subsubsection{Overfitting and ridge regression}
Increasing model complexity always lowers the training error but eventually
raises the test error \textrightarrow{} \textbf{overfitting}. A standard remedy
is to add a regularizer penalizing the weight norm:
\[
  \min_{w} \sum_{i=1}^{N} \left(x_i^T w - y_i\right)^2 + \lambda\, w^T w .
\]
For $\lambda > 0$ this stays convex. The multi-output version uses the squared
Frobenius norm:
\[
  \min_{W} \sum_{i=1}^{N} \|W^T x_i - y_i\|^2 + \lambda \|W\|_F^2 .
\]

\subsection{Linear classification}
\subsubsection{Classification as regression, and its failure}
A binary label can be predicted with a line, $\hat{y} = w^{(0)} + w^{(1)} x$, and
a sample assigned by thresholding. Multiple classes are handled with one-hot
labels and a multi-output least-square fit, assigning
$k = \arg\max_j \hat{y}^{(j)}$. This \textbf{least-square classification} fails:
a single far-away sample produces a huge residual that drags the line and
misclassifies easy points. The cause is fitting a line to inherently nonlinear
(0/1) data.

\subsubsection{Logistic regression}
We would like a discrete output; we approximate it by squashing the linear score
through a \textbf{sigmoid},
\[
  \hat{y}(x) = \sigma(w^T x) = \frac{1}{1 + \exp(-w^T x)} .
\]
For $C$ classes we use the \textbf{softmax}, giving a probability per class,
\[
  \hat{y}^{(k)} = \frac{\exp(w_{(k)}^T x)}{\sum_{j=1}^{C} \exp(w_{(j)}^T x)} ,
\]
and train by minimizing the \textbf{cross-entropy} empirical risk
\[
  R(W) = -\sum_{i=1}^{N} \sum_{k=1}^{C} y_i^{(k)} \ln \hat{y}_i^{(k)} ,
\]
via gradient descent (no closed form).

\subsubsection{Margin and the SVM}
Different classifiers give different decision boundaries, so what makes one
good? The \textbf{margin} is the orthogonal distance between the boundary and the
nearest sample; a good boundary maximizes it. The \textbf{Support Vector
Machine} is a binary classifier on the same linear model $\hat{y} = w^T x$ that
maximizes the margin while enforcing correct classification. It is written as
the convex problem
\[
  \min_{w^{(0)},\,\tilde{w}} \;
  \tfrac{1}{2}C\|\tilde{w}\|^2
  + \sum_{i=1}^{N} \max\!\left(0,\, 1 - y_i\,(w^{(0)} + \tilde{w}^T x_i)\right) ,
\]
where the $\max(\cdot)$ term is the \textbf{hinge loss} (encourages correct
classification) and $\|\tilde{w}\|^2$ is the inverse margin (minimizing it
maximizes the margin).

\subsubsection{The key insight}
\begin{definition}[Same model, different loss]
  Least-square classification, logistic regression and the SVM all rely on the
  \textit{same} linear model. What distinguishes them is only the
  \textbf{loss function}: square, logistic and hinge respectively, each
  approximating the ideal but non-differentiable $0$--$1$ loss. The choice of
  loss is what shapes the resulting algorithm.
\end{definition}

\subsection{Dealing with nonlinear data}
\subsubsection{Feature expansion}
Linear models fail on nonlinear data, but a fixed nonlinear \textbf{feature map}
$\phi(x)$ can make it separable (e.g.\ $x \to [x, x^2]$). The polynomial
expansion stacks the constant, linear, quadratic, \dots, $M$-th power terms and
cross-products, and the expanded features are then used in any algorithm,
regression or classification. The drawback: one must pick the degree (and which
functions), the number of choices grows combinatorially, and high degrees
overfit.

\subsubsection{Kernels}
Many algorithms ultimately depend only on dot products $x_i^T x_j$, which encode
similarity ($x_i^T x_j \propto \cos\angle(x_i,x_j)$). After expansion they become
$\phi(x_i)^T \phi(x_j)$. Instead of choosing $\phi$ explicitly, we directly
define a \textbf{kernel function}
\[
  k(x_i, x_j) = \phi(x_i)^T \phi(x_j) ,
\]
which corresponds to \textit{some} mapping $\phi(\cdot)$ that may be unknown.

\subsubsection{Representer theorem and kernel ridge regression}
The \textbf{representer theorem} states that the minimizer of a regularized
empirical risk is a linear combination of the training points in feature space,
\[
  w^* = \sum_{i=1}^{N} \alpha_i \phi(x_i) = \Phi^T \alpha .
\]
Substituting into ridge regression rewrites the risk purely in terms of the
kernel,
\[
  R(\alpha) = \sum_{i=1}^{N} \left(\alpha^T k(X, x_i) - y_i\right)^2
  + \lambda\, \alpha^T K \alpha ,
\]
and prediction for a new sample never needs $\Phi$ explicitly:
\[
  \hat{y} = y^T (K + \lambda I_N)^{-1} k(X, x) .
\]
The same \textbf{kernel trick} kernelizes the SVM (via its dual), logistic
regression, and later PCA. Kernel methods handle nonlinear data with often nice
(closed-form, convex) solutions, but depend on the choice of kernel and scale
with the number of samples.

\subsection{Neural networks}
Rather than design $\phi$ by hand, neural networks \textbf{learn the
representation jointly} with the final model. A layer applies a linear map
followed by an elementwise nonlinear \textbf{activation} $f$ (sigmoid, tanh,
ReLU). With the convention that $a$ is the \textbf{pre-activation} and $z$ the
\textbf{post-activation}, one hidden layer gives
\[
  z_{(1)} = f\!\left(W_{(1)}^T x\right), \qquad
  \hat{y} = f_{(2)}\!\left(W_{(2)}^T z_{(1)}\right) .
\]
Stacking layers yields a \textbf{multilayer perceptron (MLP)}; each layer has its
own parameters. Training minimizes the square loss (regression) or cross-entropy
(classification) by \textbf{stochastic gradient descent}, with gradients
computed efficiently by \textbf{backpropagation}.

\subsection{Convolutional neural networks}
Treating an image as one big vector is wasteful: it gives a huge parameter
count, ignores that nearby pixels share structure, and is not translation
robust. \textbf{Convolutions} fix this by reducing parameters and adding
translation robustness. A \textbf{CNN} stacks convolution and pooling blocks
(varying the number of channels and filter sizes), then vectorizes the final
feature map and appends fully-connected layer(s) to output a vector (class
probabilities, joint coordinates, \dots). The weights are learned exactly as in
an MLP.

\subsection{Sequences and text}
Words are mapped to dense \textbf{word embeddings}, real-valued vectors whose
similarity reflects similarity of meaning. To model a whole sentence as a
sequence:
\begin{itemize}
  \item \textbf{RNNs} process the sequence step by step, carrying a hidden state
  $h_t$ via feedback loops; unrolling them across time gives the full
  computation graph and lets them learn from history of any length.
  \item \textbf{Attention} lets the decoder focus on the relevant encoder
  outputs at each step instead of a single fixed summary vector.
  \item \textbf{Transformers} drop recurrence entirely and encode sequences with
  \textbf{multi-head self-attention} plus feedforward layers, in cascaded
  encoder and decoder blocks.
  \item The same idea returns to vision in the \textbf{Vision Transformer}, which
  treats image patches as a sequence of tokens.
\end{itemize}

\subsection{k-Nearest Neighbors}
A simple nonparametric baseline: compute the distance from the test sample $x$
to all training samples, find the nearest $x_{NN}$, and predict its label
$\hat{y} = y_{NN}$ (or a vote/average over the $k$ nearest). No training, yet
very competitive on tasks such as MNIST.

\section{Unsupervised learning}
Unsupervised learning is a single stage: analyze the whole dataset, or transform
it for further analysis, without labels.

\subsection{Clustering}
\subsubsection{K-means}
Identify $K$ groups so that within-cluster distances are small and across-cluster
distances are large, encoded via distances to cluster centers
$\{\mu_1, \dots, \mu_K\}$. The algorithm alternates until convergence:
\begin{enumerate}
  \item assign each point to the nearest center;
  \item update each center as the mean of its assigned points.
\end{enumerate}

\subsubsection{Compactness vs.\ connectivity}
K-means exploits \textbf{compactness} and fails on non-compact shapes (e.g.\
concentric rings). There, clusters are better defined by \textbf{connectivity}.

\subsubsection{Spectral clustering}
Encode the data as a graph with a similarity matrix $W \in \mathbb{R}^{N\times N}$
and a diagonal degree matrix $D$ with $D_{ii} = \sum_j W_{ij}$. Strong edges
should stay together, weak ones can be cut. The \textbf{normalized cut} is
relaxed to a generalized eigenvalue problem
\[
  (D - W)\,y = \lambda D y ,
\]
and the partition is read off the eigenvector with the \textbf{second smallest}
eigenvalue.

\subsection{Dimensionality reduction}
\subsubsection{PCA}
Real data lives on a low-dimensional manifold inside a huge ambient space. PCA
projects
\[
  y_i = W^T(x_i - \bar{x}), \qquad \text{s.t. } W^T W = I_d ,
\]
keeping the ``important'' signal by finding directions of \textbf{large
variance}. For a 1D output,
\[
  \max_{w_{(1)}} w_{(1)}^T C\, w_{(1)} \quad \text{s.t. } w_{(1)}^T w_{(1)} = 1 ,
\]
solved by the eigenvector of the covariance $C$ with the largest eigenvalue;
more dimensions take the next-largest eigenvectors. PCA also gives a mapping
\textbf{back} to the high-dimensional space, $\hat{x} = \bar{x} + W y$, which
constrains $\hat{x}$ to a subspace and thus regularizes.

\subsubsection{Kernel PCA}
Applying the kernel trick to PCA leads to the eigenvalue problem
$K a = \lambda_1 N a$, and the low-dimensional projections depend only on the
kernel,
\[
  y_i = \sum_{j=1}^{N} a_j\, k(x_i, x_j) ,
\]
so $w_{(1)}$ (which depends on the unknown $\phi$) is never needed explicitly.

\subsubsection{Fisher LDA}
PCA is unsupervised and may discard class information. Using labels, Fisher
\textbf{LDA} seeks a projection where samples of the same class cluster and
different classes separate (the same intuition as K-means, but supervised). For
a 1D projection,
\[
  \max_{w_{(1)}} w_{(1)}^T S_B\, w_{(1)} \quad \text{s.t. } w_{(1)}^T S_W\, w_{(1)} = 1 ,
\]
with $S_B$ the between-class and $S_W$ the within-class scatter matrices; the
solution is again a generalized eigenvalue problem.

\subsubsection{Autoencoders}
PCA's encode/decode pair ($y = W^T(x - \bar{x})$, $\hat{x} = \bar{x} + W y$)
generalizes to an \textbf{autoencoder}: there is no reason to keep the encoder
and decoder linear or single-layer. Stacking nonlinear layers gives a
\textbf{deep autoencoder}, and using transposed (upsampling) convolutions gives a
\textbf{convolutional autoencoder}. It is trained by minimizing the
reconstruction error
\[
  R(\{W_{(j)}\}) = \frac{1}{N} \sum_{i=1}^{N} \|\hat{x}_i(\{W_{(j)}\}) - x_i\|^2 ,
\]
non-convex and optimized by SGD, just like supervised deep networks.

\section{Cross-cutting concepts}
\subsection{Data representation}
Real data is often large and heterogeneous (e.g.\ documents of different
lengths). A good representation makes it compact and homogeneous, e.g.\
\textbf{bag-of-words} histograms of fixed length. Even then, data usually needs
pre-processing: filling in missing values, cleaning noise, and \textbf{normalizing}
features that have incommensurate magnitudes.

\subsection{Cross-validation}
Most algorithms have hyper-parameters; tuning them on the training loss favors
complex models and overfits, yet the test set must not be touched. \textbf{Cross-
validation} estimates generalization by holding out part of the training data.
The trade-offs:
\begin{itemize}
  \item \textbf{Validation set}: cheap, but an unreliable estimate.
  \item \textbf{$k$-fold} (e.g.\ 10-fold): only ``wastes'' $1/k$ of the data and
  is $k$ times the cost of a single validation set; fewer folds (3-fold) are
  cheaper but noisier.
  \item \textbf{Leave-one-out} ($N$-fold): wastes no data but is expensive.
\end{itemize}
There is no ideal solution \textrightarrow{} the choice is a cost/reliability
trade-off.
